\section{Normalizations and quantitative consequences}
\label{h:dv:applications}

The common recurrence separates choices that need not be coupled:
the local exponent, the exact-label passage, the common support,
and the use of actual or padded blocks. This section records their
substitutions and extracts the sample bounds once.
Throughout, \(d\ge16\), \(n=d-1\), \(p=\lfloor d/8\rfloor\), and
\(b=\lceil T/p\rceil\). The nearby integer \(p\) need not equal the
moment order \(q\) used in the main proof. All success probabilities
average the uniform signal, Gaussian rows, and learner randomness.
A success guarantee for every signal implies these premises by averaging.

\subsection{The density interfaces in the same recurrence}

For \(n/2\le\beta\le d/2\), a restricted continuation measure
\(\mu_w=\mathbf1_{B(z,R)}h_w\sigma_d\) with ambient coefficient \(H\)
has an exact joint density satisfying
\begin{equation}\label{h:dv:application-moment}
 \int F_w^p\,d\lambda_m
 \le C_{\mathrm D}^{d(p-1)}
       H^pR^{\beta p-m(p-1)}
 \qquad(1\le m\le p).
\end{equation}
This is Lemma~\ref{h:ev:density}; the joint-\(L^1\) or pointwise
construction gives the same bound. Either
\(\lambda_m(\mathcal E_{z,R})\le C_{\mathrm E}^dR^m\) or
\(\lambda_m(\mathcal S_{z,R})\le C_{\mathrm V}^dR^m\) restores
\(R^\beta\) in Lemma~\ref{h:alt:selection}. Thus no normalization to
a posterior and no repeated backward induction is needed.

\paragraph{Direct measurable tests.}\label{h:dv:half-ambient-application}
At \(\beta=d/2,m=p\), the inverse-distance bound for an affine
\(j\)-plane, \(0\le j\le p-2\), has the explicit form
\[
 \int\operatorname{dist}(s,L)^{-p}\,d\mu(s)
 \le\left(1+3^j2^{d/2}\frac{p}{d/2-j-p}\right)
       H R^{d/2-p},
 \qquad d/2-j-p\ge d/4 .
\]
Proposition~\ref{h:dv:measurable-test} therefore supplies the test
scale \(C^dHR^{d/2-p(1-1/p)}\), without using a strong density
conclusion. The full-block volume \eqref{h:ev:full-ellipsoid} and
the test form of Lemma~\ref{h:alt:selection} give the same coefficient
step. The hemisphere terminal coefficient and count \((p+2)2^M\)
give \eqref{h:ev:common-success} at this exponent, and hence the
base-two bounds below for either success event.

\paragraph{Joint \(L^1\) convergence.}\label{h:dv:half-sphere-application}
At \(\beta=n/2\), Proposition~\ref{h:dv:joint-l1} applies for every
actual row count \(1\le m\le p\), since \(\mu_w\ll\sigma_d\) and
\(m<d\). With the polar terminal coefficient, either support estimate
and the common recurrence give
\begin{equation}\label{h:dv:l1-success}
 \Pr_\sigma\{\operatorname{angle}(\widehat S,S)\le\epsilon\}
 \le C_0^d\epsilon^{n/2}
 \left[C_{\mathrm{L1}}^d\bigl((p+2)2^M\bigr)^{1/p}\right]^b
\end{equation}
for an absolute \(C_{\mathrm{L1}}\). Neither the quarter-net proof
nor a particular polar denominator interval is required for this
density passage.

\paragraph{Factored accuracy.}\label{h:dv:factored-application}
Writing \(H=D\epsilon^{n-\beta}\) moves the fixed accuracy factor
outside \eqref{h:ev:coefficient-recursion}; the recurrence simply
multiplies \(D\) by \(C^dN^{1/p}\).
For example, at \(\beta=n/2,m=p\), the specified pointwise density
of Proposition~\ref{h:dv:pointwise-density} satisfies
\begin{equation}\label{h:dv:factored-moment}
 \int(F_w^\sharp)^p\,d\lambda_p
 \le
 \bigl(D\epsilon^{n/2}R^{n/2}\bigr)^p
 \bigl(C_3^dR^{-p}\bigr)^{p-1}.
\end{equation}
This follows by substituting \(H=D\epsilon^{n/2}\) in
\eqref{h:dv:pointwise-moment}. The corresponding inverse-distance
bound is \(C_3^dD\epsilon^{n/2}R^{n/2-p}\), with margin at least
\(d/4\). Either net volume or the ellipsoid volume cancels the
negative radius power. The valid looser category count
\(\mathcal N=2(p+1)2^M\ge(p+2)2^M\) consequently gives
\[
 \Pr_\sigma\{\operatorname{angle}(\widehat S,S)\le\epsilon\}
 \le C_0^d\bigl(C^d\mathcal N^{1/p}\bigr)^b\epsilon^{n/2}.
\]
This enlargement has no connection to the pointwise fiber quantifier;
all category counts concern only local stopping offsets, not \(T\).

\subsection{One logarithmic extraction}

Suppose an upper bound obtained above has the form
\[
 P_{\mathrm{succ}}\le K\epsilon^\kappa e^{d b\Lambda},
 \qquad K>0,\quad\kappa,\Lambda>0 .
\]
If \(P_{\mathrm{succ}}\ge2/3\), taking natural logarithms and using
\(b\le T/p+1\) gives
\begin{equation}\label{h:dv:log-extraction}
 b\ge
 \frac{\kappa\log(1/\epsilon)-\log K-\log(3/2)}{d\Lambda},
 \qquad
 \frac Tp\ge
 \frac{\kappa\log(1/\epsilon)-\log K-\log(3/2)}{d\Lambda}-1.
\end{equation}
All the exponents above satisfy \(\kappa=n-\beta\ge d/3\).
The negative constant in this inequality is the rounded-block loss.
The full-data reflection estimate handles the remaining bounded
logarithms; the following elementary consequences retain several
useful numerical endpoints.

For a learner with a deterministic sample limit \(T\), use the
fixed-seed Borel experiment of Section~\ref{h:stream:reduction}.
Pre-generate all \(T\) Gaussian rows and grant it the full data;
its stopped output law is then a kernel of those rows and labels.
Apply Lemma~\ref{h:ev:reflection} conditionally and average the original
conditional probabilities and expectations, as in the main endpoint argument.
If its angular success
probability is at least \(2/3\) and \(\epsilon\le1/10\), then
\[
 \mathbb E\langle\widehat S,S\rangle
 \ge \frac23\cos(1/10)-\frac13>\frac14.
\]
Together with \eqref{h:ev:reflection-expectation} this proves
\begin{equation}\label{h:ev:reflection-sixteenth}
 T>d/16.
\end{equation}
For chordal success, \(T\le d/8\) and \(b_0=3/4\) in
\eqref{h:ev:antipodal} give
\begin{equation}\label{h:ev:reflection-seven-twelfths}
 \Pr\{\|\widehat S-S\|\le\epsilon\}\le\frac7{12}
 \qquad(\epsilon\le1/10).
\end{equation}
Equivalently, Markov's inequality bounds
\(\Pr\{\|u\|^2\ge3/4\}\) by \(1/6\); on its complement the reflected
points have distance greater than one, and at most half of that mass
can succeed.  The choices \(b_0=1/2\) and \(T\le d/8\) give \(5/8\).
Formula \eqref{h:ev:direct-residual} also shows that \(T\le d/5\) is
incompatible with success \(2/3\) when \(\epsilon\le1/10\).

\subsection{Explicit constants and thresholds}

All statements here assume \(0<\epsilon\le1/10\) and success at
least \(2/3\), with angular success sufficient whenever the displayed
upper bound uses chordal success. Constants may be enlarged to cover
any of the preceding density or support calculations.

\paragraph{Half the spherical dimension.}
For \(M\le d^2\), the count satisfies
\[
 \bigl((p+2)2^M\bigr)^{1/p}\le2^{32d},
\]
since \(p+2\le2^p\) and \(M/p+1\le16d+1\le32d\).
Choose \(D\ge\max\{4,e\}\) so that the per-block factor in
\eqref{h:ev:half-sphere-success} is at most \(D^d\). Then
\begin{equation}\label{h:ev:half-sphere-absolute}
 \Pr_\sigma\{\operatorname{angle}(\widehat S,S)\le\epsilon\}
 \le D^{d(b+1)}\epsilon^{n/2}.
\end{equation}
Set \(L=\log(1/\epsilon)\) and \(c_0=(3\log D)^{-1}\).
Substitution of \(K=D^d,\Lambda=\log D,\kappa=n/2\) into
\eqref{h:dv:log-extraction} gives
\(b\ge c_0L-2\) and \(T\ge p(c_0L-3)\). Hence
\begin{equation}\label{h:ev:half-sphere-threshold}
 L\ge6/c_0\quad\Longrightarrow\quad
 T\ge\frac d{16}\frac{c_0}{2}L.
\end{equation}
For \(L<6/c_0\), \eqref{h:ev:reflection-sixteenth} gives
\(T>(c_0/96)dL\).

Equivalently, absorb the initial and block factors into
\(e^{C_*d}\), \(C_*\ge\max\{1,\log4\}\). The substitution
\(K=e^{C_*d},\Lambda=C_*,\kappa=n/2\) gives
\begin{equation}\label{h:ev:exponential-threshold}
 \frac Tp\ge\frac{L}{3C_*}-3,\qquad
 L\ge18C_*\quad\Longrightarrow\quad T\ge\frac{dL}{96C_*}.
\end{equation}
This also applies to \eqref{h:dv:l1-success} when \(M=o(d^2)\):
its logarithmic category cost is \(o(d)\), so an absolute \(C_*\)
works after a dimension threshold depending on that memory sequence.
The direct residual estimate rules out \(T\le d/5\) for the bounded
complementary range.

\paragraph{Half the ambient dimension.}
For \(M=o(d^2)\), \(\log((2p+4)2^M)/p=o(d)\).
Thus an absolute \(A\ge e\), after an eventual dimension threshold
depending on the memory sequence, gives
\begin{equation}\label{h:ev:half-ambient-absolute}
 \Pr_\sigma\{\|\widehat S-S\|\le\epsilon\}
 \le A^{d(b+1)}\epsilon^{d/2-1}.
\end{equation}
With \(a_0=(3\log A)^{-1}\), substitute
\(K=A^d,\Lambda=\log A,\kappa=d/2-1\) in
\eqref{h:dv:log-extraction} to obtain
\begin{equation}\label{h:ev:half-ambient-log}
 b\ge a_0L-2,\qquad T\ge p(a_0L-3),
\end{equation}
and therefore
\begin{equation}\label{h:ev:half-ambient-threshold}
 L\ge6/a_0\quad\Longrightarrow\quad
 T\ge\frac d{16}\frac{a_0}{2}L.
\end{equation}
For \(L<6/a_0\), \eqref{h:ev:reflection-seven-twelfths} forces
\(T>d/8>(a_0/48)dL\).

For the base-two form, use instead the hemisphere coefficient
\(4^{d-1}\epsilon^{d/2-1}\) and the count \((p+2)2^M\).
When \(M\le d^2\), an absolute \(C_1\ge1\) bounds the block factor
by \(2^{C_1d}\). With \(L_2=\log_2(1/\epsilon)\), the same
logarithmic extraction, using
\(2(d-1)-\log_2(2/3)\le3d\), gives
\begin{equation}\label{h:ev:binary-endgame}
 \Pr_\sigma\{\|\widehat S-S\|\le\epsilon\}
 \le2^{C_1db}4^{d-1}\epsilon^{d/2-1},
 \qquad
 \frac Tp\ge\frac{L_2/3-3}{C_1}-1.
\end{equation}
In particular,
\begin{equation}\label{h:ev:binary-threshold}
 L_2\ge6(3+C_1)\quad\Longrightarrow\quad
 T\ge\frac{dL_2}{96C_1}.
\end{equation}
Here \(L_2/3-3-C_1\ge L_2/6\) and \(p\ge d/16\).
In the complementary range, the antipodal choice \(b_0=1/2\)
gives \(5/8<2/3\) when \(T\le d/8\).
These formulas also apply to the direct-test route above.

\paragraph{The accuracy-factored form.}
For \(M=o(d^2)\), the count in the factored recurrence obeys
\(\mathcal N^{1/p}\le e^d\) eventually. For absolute constants
\(C_4,C_5\) it follows that
\begin{equation}\label{h:dv:factored-endpoint}
 \Pr_\sigma\{\operatorname{angle}(\widehat S,S)\le\epsilon\}
 \le e^{C_4d(b+1)}\epsilon^{(d-1)/2},
 \qquad
 \log(1/\epsilon)\le C_5(T/p+2).
\end{equation}
The second inequality is \eqref{h:dv:log-extraction}, with
\(n/2\ge d/3\) and the fixed \(\log(3/2)\) absorbed into \(C_5\).
Reflection gives \(T>d/8\), and \(p\ge d/16\) then gives
\(T/p+2<32T/d\). Consequently
\(T\ge d\log(1/\epsilon)/(32C_5)\) over the entire accuracy range,
without a relation between precision and memory.

For each fixed \(A_0>0\), replacing the memory condition by
\(M\le A_0d^2\) changes only the constants and bounded-logarithm
thresholds, which may depend on \(A_0\). These arguments do not give
a uniform positive constant for unbounded \(A_0=A_0(d)\).
